\chapter{Veranda Calc}
\label{app:veranda}

\textbf{Veranda Calc} is \ccprod{} on Veranda, the GEOS-style desktop for the
Commander X16. The worksheet, the formulas, the files and the keys are the
ones this manual describes; what follows is what is different.

\section{Installing and starting}

Veranda Calc is one file, \texttt{Veranda Calc.cvt}, with all its overlays
inside it. Copy it to the card beside Veranda (the top folder is the best
place: see below) and double-click its icon in Veranda's shell.

A double click on a workbook --- a file ending \texttt{.X16S},
\texttt{.XLSX}, \texttt{.ODS} or \texttt{.CSV} --- also starts Veranda Calc,
with that file open: an \texttt{.X16S} workbook is opened, the others are
imported (Chapter~\ref{ch:importexport}). The shell looks for Veranda Calc
in the workbook's own folder first, then in the top folder. Long names
work: a workbook called \texttt{Household budget for 2026.X16S} opens
with a double click like any other. \menupath{File > Quit} goes back to
the shell, in the folder it was showing. In the shell, Veranda Calc shows
as its own tile, a green sheet of cells.

\section{The screen}

Veranda Calc looks like Veranda's own programs (Veranda's ``Look~2''). At
the top is Veranda's menu bar, with the Veranda mark, the clock and the
sound; under it the \textbf{title band}, which names the workbook
(\textit{Untitled} until it has been saved or opened). The formula bar is
the toolbar: the cell's name on grey, what is in it in a white field.
Then the grid --- the same eighty columns of character cells, its row and
column headers grey --- and at the bottom the sheets' tabs (the one shown
in white, joined to its sheet) and the \textbf{status bar}: \textsf{READY},
\textsf{EDIT} or \textsf{MODIF}, the cell, and the memory free. The
cursor, a selected block and the headers of the cursor's row and column
are in the selection colour, light blue. The menu bar and the title band
take four card rows, so the grid has three rows fewer than on the plain
X16. Cell borders are drawn as lines on the edges of the cells. Dark mode
(Veranda's Preferences) gives the same screen on Veranda's dark palette.

\section{Menus and dialog boxes}

The menus are Veranda's: click a title, then an item. The first,
\menupath{Veranda Calc}, has \menupath{About Veranda Calc}, which names
the version, and \menupath{Desk accessory}, which runs one of the desk
accessories installed with the applications (Preferences, the
Calculator, the Note Pad) over the workbook. Then \menupath{File},
\menupath{Edit}, \menupath{Layout}, \menupath{Data}, \menupath{Chart} and
\menupath{Sheet}, as on the plain X16; the keys shown beside the items
work as before. Every question --- Find, Replace, a column width, a
sheet's name, ``discard unsaved changes?'' --- is a Veranda dialog box with
buttons; \key{RETURN} answers OK.

\section{The pointer}

A click moves the cursor, a drag selects a block, the wheel scrolls, and a
click on a sheet's tab shows that sheet. The \textbf{right button} edits
the cell under the pointer, as \key{F2} does.

\section{Opening, importing and saving}

\menupath{File > Open}, \menupath{File > Import}, \menupath{File > Save
as} and \menupath{File > Export} use Veranda's standard Open and Save
box, the one Veranda's own programs use. Its title band can be dragged
to move it. It lists the folders first, each with a yellow folder badge,
then the files, each with a badge of its kind; names of up to 62
characters are shown, long ones cut with ``\ldots''; hidden names
(starting with a dot) are not. Up and down, \key{PgUp}, \key{PgDn},
\key{Home} and \key{End} move the selection, the wheel moves it three at a
time; \key{RETURN}, the \textbf{Open} button or a second click on the
selected name opens it. Open on a folder goes into it, and \texttt{..}
goes up one, as far as the top folder of the card; when opening, a letter
goes to the next name starting with it and \key{Backspace} goes up a
folder. \textbf{In:} shows the folder; a click on it lists the folders
above it and Veranda's places --- Documents, Apps and the SD card --- to go
to. \textbf{New folder} makes a folder (type its name, \key{RETURN}) and
goes into it. The type button switches between the files Veranda Calc
reads (\textsf{Workbooks}, \texttt{.X16S}, for Open and Save; \textsf{Sheets},
\texttt{.XLSX}, \texttt{.ODS} and \texttt{.CSV}, for Import
and Export) and \textsf{All files}. \key{ESC} or \textbf{Cancel} gives
up.

The box starts in the folder Veranda Calc is working in: the shell's
folder when it starts, then the folder of the last file opened, imported
or saved. Save writes a workbook back into the folder it came from, under
its own name, however long. Save as and Export start with the workbook's
name in the field and save into the folder the box shows; Save as adds
\texttt{.X16S} to a name typed without an extension. A file's name and its
folders together may be up to 63 characters long.

A workbook that has just been opened is not ``modified'': New and Quit ask
nothing until it is changed.

\section{Charts}

\menupath{Chart} draws bars, a line or a pie in black and white: bars and
slices are told apart by their fill patterns, and the pie's key shows each
slice's pattern beside its label. Any key or a click puts the sheet back;
\key{S} first writes the chart to \texttt{CHART.BMP}, a 640 by 464
black-and-white picture.

\section{Memory and speed}

Veranda keeps some of the X16's banked memory for itself and for a desk
accessory, so a 512\,KB machine has about 320\,KB for workbooks. Opening a
large workbook is no faster than on the plain X16: a sheet of 2,400 cells
takes about forty seconds to open and thirty to save.
